% \begin{figure}[h!]
% \centering
% \begin{forest}
% for tree={
%     draw,
%     rounded corners,
%     align=center,
%     minimum width=2.2cm,
%     minimum height=1cm,
%     inner sep=2pt,
%     font=\small
% },
% good/.style={fill=green!20},
% bad/.style={fill=red!20},
% neutral/.style={fill=blue!10}
% [P$_0$\\(root)\\{\color{blue}select $x_1$}, neutral
%     [P$_1$\\{\color{blue}$x_1 \le 0$}, good,
%         edge label={node[midway,above,font=\scriptsize]{branch}}
%     ]
%     [P$_2$\\{\color{blue}$x_1 \ge 1$}\\{\color{blue}select $x_2$}, neutral,
%         edge label={node[midway,above,font=\scriptsize]{branch}}
%         [P$_5$\\{\color{blue}$x_2 \le 0$}, bad,
%             edge label={node[midway,above,font=\scriptsize]{prune}}
%         ]
%         [P$_6$\\{\color{blue}$x_2 \ge 1$}, good,
%             edge label={node[midway,above,font=\scriptsize]{feasible}}
%         ]
%     ]
% ]
% \end{forest}
% \caption{Illustration of Branch-and-Bound showing variable selection, branching, and bounding/pruning.}
% \label{fig:branch_and_bound_tree}
% \end{figure}


\begin{figure}[htbp]
\centering

\definecolor{bbActive}{RGB}{225,225,248}
\definecolor{bbFeasible}{RGB}{210,245,210}
\definecolor{bbPruned}{RGB}{255,215,215}
\definecolor{bbBlue}{RGB}{40,80,200}

\tikzset{
    branchlabel/.style={
        midway,
        fill=white,
        inner sep=1pt,
        font=\scriptsize,
        text=black!75
    }
}

\begin{forest}
for tree={
    draw=black!75,
    rounded corners=3pt,
    line width=0.45pt,
    align=center,
    font=\small,
    inner sep=3pt,
    minimum width=2.1cm,
    minimum height=0.85cm,
    edge={-latex, line width=0.45pt, draw=black!70},
    l sep=9mm,
    s sep=10mm,
},
active/.style={fill=bbActive},
feasible/.style={fill=bbFeasible},
pruned/.style={fill=bbPruned}
[
    {$P_0$\\[-1pt]\footnotesize root\\[-1pt]\textcolor{bbBlue}{select $x_1$}},
    active
    [
        {$P_1$\\[-1pt]\textcolor{bbBlue}{$x_1 \le 0$}},
        feasible,
        edge label={node[branchlabel, above left]{branch}}
    ]
    [
        {$P_2$\\[-1pt]\textcolor{bbBlue}{$x_1 \ge 1$}\\[-1pt]\textcolor{bbBlue}{select $x_2$}},
        active,
        edge label={node[branchlabel, above right]{branch}}
        [
            {$P_5$\\[-1pt]\textcolor{bbBlue}{$x_2 \le 0$}},
            pruned,
            edge label={node[branchlabel, above left]{prune}}
        ]
        [
            {$P_6$\\[-1pt]\textcolor{bbBlue}{$x_2 \ge 1$}},
            feasible,
            edge label={node[branchlabel, above right]{feasible}}
        ]
    ]
]
\end{forest}

\caption{Illustration of a branch-and-bound search tree. At each branching node, a variable is selected and the current subproblem is split into child subproblems. Nodes can then be pruned by bounding or retained when they yield feasible incumbent solutions.}
\label{fig:branch_and_bound_tree}
\end{figure}